\section{Integration, reproducibility, and provenance}\label{sec:integration}

\subsection{Package layout and proposed repository placement}

The archive contains the complete modified \code{fast/} tree and an
untouched \code{reference/fast/} snapshot from the pinned revision.
The integration patch is relative to
\code{Topology/UnknotRecognition/fast}. It adds the modules and tests
described below and the optional source-profile argument; it does not
alter the default recognition dispatch.
The article and research artifacts can be placed in a new report
directory under \code{Topology/UnknotRecognition/reports/}, using the
next free report number at integration time.

\begin{center}
\small
\begin{tabular}{p{0.43\linewidth}p{0.48\linewidth}}
\toprule
File or directory & Purpose\\
\midrule
\code{article/main.tex} and section files &
Complete article source, with an ordinary LaTeX bibliography.\\
\code{article/main.pdf} & Compiled and visually checked article.\\
\code{article/figures/} & Vector figures and PNG previews, generated
from exact output or retained measurements.\\
\code{fast/fastunknot/twist/tail.py} &
Exact one-reference homological recurrence and compact-profile API.\\
\code{fast/fastunknot/twist/streaming.py} &
Bounded-cache, degree-streamed full homology and early recognition.\\
\code{fast/fastunknot/braid_profile.py} &
Run-encoded structural certificates and the dominated signature
experiment.\\
\code{fast/fastunknot/twist/continuation.py} &
Raw RLE command-line and Python frontend.\\
\code{fast/tests/} & All inherited and new regression tests.\\
\code{tools/} & Portable audits, benchmarks, figure regeneration,
validation, and article build scripts.\\
\code{results/} & Raw measurements, pilot runs, exact comparison
records, and the complete test log.\\
\code{reference/} & Pinned baseline and detailed proof/accounting
audit notes.\\
\code{integration.patch} & Additive patch for review against the
pinned baseline.\\
\bottomrule
\end{tabular}
\end{center}

The baseline MIT No Attribution (MIT-0) license is retained.
No remote repository write or
default-backend replacement is part of this deliverable.
The accompanying integration guide gives the exact paths and patch
check command.

\subsection{Interfaces and examples}

From the package's \code{fast/} directory, the new frontend accepts
run-encoded JSON:
\begin{lstlisting}
{"strands": 3, "runs": [[1, 1001], [2, -1], [1, 1], [2, -1]]}
\end{lstlisting}
The corresponding commands are:
\begin{lstlisting}
python -m fastunknot.twist.continuation input.json
python -m fastunknot.twist.continuation input.json \
  --mode homology --method tail --check-d2
python -m fastunknot.twist.continuation input.json \
  --mode homology --method streaming
python -m fastunknot.twist.continuation input.json --mode profile
\end{lstlisting}

The default is recognition with the tail method selected for any
necessary exact endpoint. Recognition first checks the original
one-component closure and tries the direct structural certificate.
Only an inconclusive certificate reaches the selected backend.
Full homology mode always computes the selected full exact homology,
including for links; it does not stop after a cheap knot verdict.
Profile mode returns only the structural certificate.
The methods \code{tail}, \code{streaming}, and \code{macro} are explicit
choices.

A Python call avoids JSON serialization limits for extremely large
integers:
\begin{lstlisting}
from fastunknot.twist.core import Run
from fastunknot.twist.tail import tail_homology, profile_dimension

m = (1 << 100000) + 1
answer = tail_homology(2, [Run(1, m)], check_d2=True)
assert answer["reduced_rank"] == m
assert answer["reference"]["stats"]["basis"] == 4
assert profile_dimension(answer["degree_profile"], m) == 1
\end{lstlisting}

The standard Python JSON integer conversion has a decimal-digit
limit in this runtime. The API supports the tested 100,001-bit
integers directly; the CLI does not disable that general runtime
limit. JSON also converts dictionary degree keys to strings.
The documentation identifies the schema and inclusive interval
semantics so downstream consumers do not confuse a large encoded
endpoint with an explicit array index range.

Existing checked diagrams can opt into the source specialization via
\code{recognize(diagram, use_braid_profile=True)} or the production
\code{--braid-profile} option. It is used only with retained checked
braid source. A PD-only diagram does not acquire invented braid
provenance, and a later diagram factor is not assumed to share the
original source braid.

\Needspace{15\baselineskip}
\subsection{Reproduction commands}

From the archive root:
\begin{lstlisting}
PYTHONPATH=fast python -m unittest discover -s fast/tests -v
python tools/audit_long_twist.py
python tools/audit_slope.py
python tools/benchmark_tail.py
python tools/benchmark_streaming.py --rounds 7 --batch 5
python tools/benchmark_braid_profile.py --rounds 7
python tools/make_figures.py
bash tools/build_article.sh
\end{lstlisting}

The exact outputs originally used in the article are retained.
Rerunning measurements overwrites their default target files; retain a
copy or specify a separate output path for comparison.
Wall-clock measurements on another machine are expected to differ.
The scripts use package-relative imports, fixed stated random seeds,
and standard-library exact arithmetic for the core algorithms.
Matplotlib is used only to regenerate plots.
Building the article requires a conventional LaTeX installation with
the packages named in the preamble and \code{latexmk}.

The independent recurrence audit was executed successfully before
being relocated into the final package. The relocated script changes
paths and explanatory text, not the tested algebra, and its raw result
file is retained unchanged. Its provenance note states this explicitly.
The final integrated 254-method suite was run on the delivered combined
implementation.

\subsection{Evidence boundaries and integration decisions}

The proofs establish a theorem for finite complexes obtained from the
stated local normal form. The implementation comparisons check that
the delivered code computes that model on the tested inputs.
Neither is a proof-assistant formalization of the entire topology
dependency chain.
Benchmarks establish observations on specified input families and
execution conditions; they do not establish a distribution-free
performance claim.

For integration, the evidence supports adding the tested APIs, research
documentation, and opt-in source-profile stage. The streamed backend
should remain selectable on memory grounds while its measured time
loss is addressed. A future automatic dispatcher should be justified
on the subset of inputs reaching the expensive endpoint, with both
time and memory objectives stated.

The main unresolved theorem is universal closure of a compact
representation under construction, gluing, reduction, and exact
decision. The finite-threshold result supplies one such closure
mechanism in a controlled setting, together with enough source,
proof, and experimental detail to make the next stage of research
reviewable and reproducible.
